\chapter{Formát PDF/A}

Трансформация понятия «мышление»: от человеческого к вычислительному За историей технических проектов искусственного разума стоит глубокая концептуальная трансформация: само понятие «мышления» последовательно сужалось и переопределялось так, чтобы машина могла в него вписаться.

В античности и в значительной мере в Средневековье мышление понималось как нечто несводимое к логическому выводу: оно включало нравственное суждение, эстетическое восприятие, мудрость как способность различать добро и зло в конкретной ситуации.
Aristotel'evskij nous был прежде всего способностью к созерцанию высших истин, а не механизмом дедукции.
Сенека в «Нравственных письмах» настаивал на том, что разум человека «причастен нравственным категориям, которые уходят гораздо дальше, чем просто логическое суждение — добро, зло, красота, уродство».
Именно поэтому «то, что действительно отличает человека от животного это способность к сложнейшей операции соединения логического и нравственного, нравственному выбору как таковому».
Это отличительная черта человеческого сознания: «выбирать не только рационально, но и нравственно».
Картезианский механицизм произвёл первое принципиальное сужение: мышление стало отождествляться с рациональным, логическим компонентом — тем, что потенциально поддаётся формализации.
Нравственное, эмоциональное, эстетическое было выведено за скобки или атрофировалось к «страстям», подлежащим управлению.
Лейбницевский проект «исчисления разума» сделал следующий шаг: мышление стало отождествляться с формальными операциями над символами.
Логика Буля, а затем математическая логика Фреге и Рассела завершили этот процесс на уровне формальных систем.
Тьюринг поставил принципиальный вопрос: если мы не можем отличить ответы машины от ответов человека в диалоге, есть ли у нас основания отрицать, что машина мыслит?
Этот вопрос был революционным — но именно потому, что он предполагал редукцию мышления к поведению, наблюдаемому извне.
Внутренний опыт, нравственная укоренённость, интенциональность — всё это было вынесено за скобки как ненаблюдаемое и потому нерелевантное.
Сёрл показал несостоятельность этой редукции: его «Китайская комната» демонстрирует, что система может производить правильные ответы, не понимая их — не имея никакого доступа к смыслу манипулируемых символов (Minds, Brains and Science, 1984).
Дрейфус добавил к этому аргумент воплощённости: даже если отвлечься от нравственного измерения, человеческое познание неотделимо от телесного опыта, от «знания-как» в отличие от «знания-что», — и никакая формальная система не воспроизведёт эту укоренённость (What Computers Can't Do, 1972).
Тем не менее прагматическая, редукционистская концепция мышления-как-вычисления оказалась исторически победоносной, но не потому, что она философски убедительнее, а потому, что она технически продуктивна и экономически выгодна.
Как показывает Петрунин, идея ИИ «опирается на господствующие в обществе интуиции» (Петрунин, б/г, с. 12) — и господствующей интуицией нашего времени оказалась интуиция технократическая: разум есть то, что можно измерить, формализовать и воспроизвести.
Всё остальное либо несущественно, либо будет формализовано позже.
Последствия этой победы выходят далеко за пределы технических дискуссий.
Если мышление — это вычисление, то алгоритм может принимать решения вместо человека: о кредитоспособности, о рисках, о виновности, о пригодности к должности.
Если интеллект измеряется поведенческим тестом, то система, проходящая тест, обладает достаточным интеллектом для управления.
Если нравственное измерение мышления нерелевантно, то вопрос об этической ответственности алгоритма не возникает.
Олоф-Орс и Смит показывают, что именно это концептуальное наследие делает возможным антропоморфный дизайн как инструмент надзорного капитализма: система, воспроизводящая поведение разумного существа, воспринимается пользователем как разумное существо и весь аппарат социального доверия, предназначенный для взаимодействия с другими людьми, оказывается задействован в интересах корпорации (Olof-Ors \& Smit, 2025).
В история от Луллия до ChatGPT мы наблюдаем последовательное переопределение мышления, чтобы машина могла в него вписаться.
На каждом этапе что-то важное исключалось из понятия: сначала нравственное суждение, потом интенциональность, потом воплощённость, наконец — само понимание.
Что осталось?
Статистически правдоподобная последовательность токенов.
Но это и не то, что философская традиция называла мышлением.
И это различие отнюдь не академическое: от него зависит, кто несёт ответственность за решения, принимаемые «интеллектуальными» системами, и чьи интересы эти решения обслуживают.

/chapter